\begin{center}
    {\LARGE \textbf{ĐỀ LUYỆN 05}}\\[6pt]
    {\large (Dành cho đội tuyển quốc gia)}
\end{center}

\begin{tabularx}{\textwidth}{|c|X|c|c|c|}
\hline
\textbf{Bài} & \textbf{Tên bài} & \textbf{File đầu vào} & \textbf{File đầu ra} & \textbf{Điểm} \\
\hline
1 & Hai đống & \texttt{TPILE.INP} & \texttt{TPILE.OUT} & 7 \\
\hline
2 & Tổ hợp thú vị & \texttt{CFUN.INP} & \texttt{CFUN.OUT} & 7 \\
\hline
3 & Dãy XOR tăng chặt & \texttt{XARR.INP} & \texttt{XARR.OUT} & 6 \\
\hline
\end{tabularx}



\section*{Bài 1: Hai đống}
% https://www.codechef.com/problems/SPLITMIN

Bạn được cho $N$ cặp số nguyên. Cặp thứ $i$ là $(A_i, B_i)$.  
Bạn có hai đống (pile), ban đầu đều rỗng.

Bạn sẽ thực hiện quy trình sau:

\begin{itemize}
    \item Với mỗi $i$ từ $1$ đến $N$, chọn \textbf{một trong hai} số $A_i$ hoặc $B_i$.
    \item Sau đó, thêm số đã chọn vào \textbf{một trong hai đống}.
\end{itemize}

Cuối quá trình, \textbf{cả hai đống đều phải không rỗng} (mỗi đống phải có ít nhất một phần tử).

Giá trị của một đống được định nghĩa là \textbf{phần tử lớn nhất} trong đống đó.

Bạn được phép chọn tối ưu để \textbf{giảm thiểu chênh lệch} giữa hai giá trị đống.

Nhiệm vụ:  
\[
\text{Tìm giá trị nhỏ nhất có thể của } \bigl| \max(\text{pile}_1) - \max(\text{pile}_2) \bigr|.
\]

\subsection*{Dữ liệu vào}

\begin{itemize}
    \item Dòng đầu tiên: số nguyên $T$ — số lượng test.
    \item Với mỗi test:
    \begin{itemize}
        \item Dòng đầu: số nguyên $N$.
        \item $N$ dòng tiếp theo: mỗi dòng chứa hai số nguyên $A_i, B_i$.
    \end{itemize}
\end{itemize}

\subsection*{Dữ liệu ra}

In ra một số nguyên trên mỗi test — giá trị nhỏ nhất có thể của hiệu tuyệt đối giữa hai đống.

\subsection*{Ràng buộc}

\begin{itemize}
    \item $1 \le T \le 10^{4}$
    \item $2 \le N \le 2 \cdot 10^{5}$
    \item $1 \le A_i, B_i \le 10^{9}$
    \item Tổng $N$ trên tất cả test không vượt quá $2 \cdot 10^5$
\end{itemize}

\subsection*{Sample Input}

\begin{verbatim}
3
2
1 6
4 9
2
1 9
4 6
4
10 10
23 55
21 80
105 91
\end{verbatim}

\subsection*{Sample Output}

\begin{verbatim}
2
3
11
\end{verbatim}

\subsection*{Giải thích}

\textbf{Test 1.}  
Các cặp là $(1, 6)$ và $(4, 9)$.  
Chọn $6$ từ cặp 1 và $4$ từ cặp 2, rồi đặt vào hai đống khác nhau.  
Khi đó:
\[
|6 - 4| = 2.
\]
Không thể đạt giá trị nhỏ hơn.

\textbf{Test 2.}  
Các cặp là $(1,9)$ và $(4,6)$.  
Chọn $1$ vào đống 1 và $4$ vào đống 2:
\[
|1 - 4| = 3.
\]

\textbf{Test 3.}  
Một nghiệm tối ưu:
\begin{itemize}
    \item Cặp 1: đặt $10$ vào đống 1.
    \item Cặp 2: đặt $55$ vào đống 2.
    \item Cặp 3: đặt $80$ vào đống 2.
    \item Cặp 4: đặt $91$ vào đống 1.
\end{itemize}
Hai đống có giá trị:
\[
\max(\text{pile}_1)=91,\quad \max(\text{pile}_2)=80.
\]
Hiệu:
\[
|91 - 80| = 11.
\]
Đây là giá trị nhỏ nhất có thể.

\newpage

\section*{Bài 2. Tổ hợp thú vị}
% https://www.codechef.com/problems/COMBI_IS_FUN

Xét một đa tập $S$ và một số nguyên dương $X$.  
Giả sử ta phân hoạch $S$ thành $L$ tập con không rỗng sao cho mỗi phần tử của $S$ thuộc đúng một tập con.

Chi phí của một phép phân hoạch được xác định bởi:
\[
\sum_{i=1}^{L} \left( X + (\max_i - \min_i) \right),
\]
trong đó $\max_i$ và $\min_i$ là giá trị lớn nhất và nhỏ nhất trong tập con thứ $i$.

Kí hiệu $F(S, X)$ là chi phí nhỏ nhất có thể đạt được nếu ta phân hoạch $S$ một cách tối ưu.

Bạn được cho các số nguyên $N, M, D, K$.  
Hãy đếm số lượng đa tập $A$ sao cho:

\begin{itemize}
    \item $A$ có đúng $N$ phần tử;
    \item $1 \le A_i \le M$ với mọi $1 \le i \le N$;
    \item $F(A, D) = K$.
\end{itemize}

Vì ta xét đa tập nên các hoán vị giống nhau chỉ được tính một lần, ví dụ  
$\{1,2,2\}$ và $\{2,2,1\}$ được xem là giống nhau.

Kết quả cần được in theo modulo $998244353$.

\subsection*{Dữ liệu vào}

\begin{itemize}
    \item Dòng đầu chứa số nguyên $T$ — số lượng test.
    \item Mỗi test gồm bốn số nguyên $N, M, D, K$.
\end{itemize}

\subsection*{Dữ liệu ra}

Với mỗi test, in ra một số nguyên — số lượng đa tập thỏa mãn điều kiện, modulo $998244353$.

\subsection*{Ràng buộc}

\begin{itemize}
    \item $1 \le T \le 10^3$
    \item $1 \le N \le 10^3$
    \item $1 \le M, D, K \le 2 \cdot 10^6$
    \item Tổng $N$ của tất cả test không vượt quá $5 \cdot 10^3$
    \item Tổng $M$ của tất cả test không vượt quá $5 \cdot 10^6$
\end{itemize}

\subsection*{Sample Input}

\begin{verbatim}
5
2 3 2 4
3 3 2 4
2 5 4 10
20 25 5 12
343 343 3 43
\end{verbatim}

\subsection*{Sample Output}

\begin{verbatim}
1
3
0
9143199
840977574
\end{verbatim}

\subsection*{Giải thích}

\textbf{Test 1.}  
Đa tập duy nhất thỏa mãn là $\{1,3\}$.  
Phân hoạch có thể là:
\[
\{1\}, \{3\} \quad \Rightarrow \quad (2 + 0) + (2 + 0) = 4,
\]
hoặc một nhóm:
\[
\{1,3\} \quad \Rightarrow \quad 2 + (3 - 1) = 4.
\]

\textbf{Test 2.}  
Các đa tập thỏa mãn:
\[
\{1,1,3\},\ \{1,3,3\},\ \{1,2,3\}.
\]
Ví dụ với $\{1,2,3\}$:
\[
2 + (3 - 1) = 4.
\]

\textbf{Test 3.}  
Không tồn tại đa tập thỏa mãn chi phí bằng $K$.

\subsection*{Subtask}

\begin{itemize}
  \item (30\% điểm) $N \le 10$, $M, D, K \le 100$.
  \item (70\% điểm) Không có ràng buộc thêm.
\end{itemize}

\newpage

\section*{Bài 3: Dãy XOR tăng chặt}

Bạn được cho một mảng $A$ gồm $N$ số nguyên dương đôi một khác nhau.

Bạn được phép thực hiện đúng \textbf{một thao tác} sau:

\begin{itemize}
    \item Chọn một số nguyên không âm $X$.
    \item Với mọi $1 \le i \le N$, cập nhật:
    \[
        A_i := A_i \oplus X,
    \]
    trong đó $\oplus$ là phép XOR theo bit.
\end{itemize}

Nhiệm vụ của bạn là tìm:

\[
\text{Độ dài lớn nhất có thể của \textbf{một đoạn con tăng chặt} trong mảng sau khi thực hiện đúng một thao tác.}
\]

\textbf{Định nghĩa:}
\begin{itemize}
    \item Đoạn con (subarray) là đoạn liên tiếp của mảng.
    \item Một đoạn con được gọi là \textbf{tăng chặt} nếu các phần tử của nó thỏa:
    \[
        a_1 < a_2 < \cdots < a_k.
    \]
\end{itemize}

\subsection*{Dữ liệu vào}

\begin{itemize}
    \item Dòng đầu: số nguyên $T$ — số lượng test.
    \item Với mỗi test:
    \begin{itemize}
        \item Dòng đầu chứa số $N$.
        \item Dòng tiếp theo chứa $N$ số nguyên đôi một khác nhau: $A_1, A_2, \ldots, A_N$.
    \end{itemize}
\end{itemize}

\subsection*{Dữ liệu ra}

Với mỗi test, in ra một số nguyên — độ dài lớn nhất của một đoạn con tăng chặt sau đúng một phép XOR với một số $X$ tùy chọn.

\subsection*{Ràng buộc}

\begin{itemize}
    \item $1 \le T \le 1000$
    \item $1 \le N \le 10^5$
    \item $1 \le A_i < 2^{30}$
    \item Các phần tử của $A$ là \textbf{phân biệt}
    \item Tổng $N$ của tất cả test không vượt quá $10^6$
\end{itemize}

\subsection*{Sample Input}

\begin{verbatim}
3
3
3 2 5
4
10 20 30 40
5
14 11 9 3 7
\end{verbatim}

\subsection*{Sample Output}

\begin{verbatim}
3
4
4
\end{verbatim}

\subsection*{Giải thích}

\textbf{Test 1.}  
Chọn $X = 1$. Khi đó:
\[
[3 \oplus 1,\ 2 \oplus 1,\ 5 \oplus 1] = [2,3,4],
\]
đoạn tăng chặt dài nhất có độ dài $3$.

\textbf{Test 2.}  
Chọn $X = 0$. Mảng giữ nguyên:
\[
[10,20,30,40],
\]
đoạn tăng chặt dài nhất có độ dài $4$.

\textbf{Test 3.}  
Chọn $X = 15$:
\[
[14\oplus 15,\, 11\oplus 15,\, 9\oplus 15,\, 3\oplus 15,\, 7\oplus 15]
= [1,4,6,12,8].
\]
Đoạn tăng chặt dài nhất là $[1,4,6,12]$ có độ dài $4$.

\subsection*{Subtask}

\begin{itemize}
  \item (30\% điểm) $N \le 2000$.
  \item (70\% điểm) Không có ràng buộc thêm.
\end{itemize}